% !TEX root = ../../../main.tex
% 所属：第一部分：单变量分析学 / 第6章：函数的极限与连续性
% 节：6.5 无穷大量与无穷小量的阶
% 正文按原文件顺序保存；由 ../chapter6.tex 调用。
\section{无穷大量与无穷小量的阶}
下面我们介绍阶数这一概念，类似于集合的势的比较，我们定义的阶数揭示了无穷小量趋近于零的相对程度以及无穷大量绝对增长至无穷的相对程度

类似于数列极限中对无穷大量与无穷小量的定义，下面是函数极限中无穷大量与无穷小量 的定义
\begin{definition}
	若 $\lim\limits_{x \to A} f(x) = 0$，则称 $x \to A$ 时 $f(x)$ 是无穷小量.
	
	若 $\lim\limits_{x \to A} f(x) = \infty(\pm\infty
	)$ ，则称 $x \to A$ 时 $f(x)$ 是无穷大量.
	
	这里 $x \to A$ 包括了 $x \to x_0$，$x \to x_0^+$，$x \to x_0^-$，$x \to +\infty$，$x \to -\infty$，$x \to \infty$ 等各种情况，下面用 $x \to x_0$ 举例.
	
\end{definition}
无穷大量与无穷小量的阶定义如下，它是一种相对的定义我们可以说两个无穷量的阶数如何如何，想要给出某一无穷量的绝对阶数，需要与$x^a$做比对
\begin{definition}
	对于 $u(x), v(x)$，当 $x \to x_0$ 时，$u(x), v(x)$ 都是无穷小量.
	
	(1) 若 $\lim\limits_{x \to x_0} \dfrac{u(x)}{v(x)} = 0$，我们称 $x \to x_0$ 时，$u(x)$ 关于 $v(x)$ 是高阶无穷小量或 $v(x)$ 关于 $u(x)$ 是低阶无穷小量，记作
	\[ u(x) = o(v(x)) \quad (x \to x_0). \]
	
	(2) 若 $\exists A > 0$，当$x$ 在$x_0$的某邻域内，成立
	\[ \left| \frac{u(x)}{v(x)} \right| \leqslant A \]
	则称 $x \to x_0$ 时，$\dfrac{u(x)}{v(x)}$ 是有界量，记作
	\[ u(x) = O(v(x)) \quad (x \to x_0). \]
	
	(3) 若 $\exists A\geqslant a > 0$，对于 $x$ 在$x_0$的某去心邻域内，成立
	\[ a \leqslant \left| \frac{u(x)}{v(x)} \right| \leqslant A \]
	则称 $x \to x_0$ 时，$u(x)$ 与 $v(x)$ 是同阶无穷小量.
	
	若 $\lim\limits_{x \to x_0} \dfrac{u(x)}{v(x)} = C \neq 0$，则 $u(x)$ 与 $v(x)$ 必是同阶无穷小量.
	
	(4) 若 $\lim\limits_{x \to x_0} \dfrac{u(x)}{v(x)} = 1$，则称 $x \to x_0$ 时，$u(x)$ 与 $v(x)$ 是等价无穷小量，记作
	\[ u(x) \sim v(x) \quad (x \to x_0). \]
	
	也可以写成
	\[ u(x) = v(x) + o(v(x)) \quad (x \to x_0). \]
	
	表示$x \to x_0$ 时，$u(x)$ 与 $v(x)$ 不一定相等，差一个关于 $v(x)$ 的高阶无穷小量，
	
\end{definition}
\begin{example}
	(1)
	\[
	\lim_{x \to 0} \frac{1 - \cos x}{x} = \lim_{x \to 0} \dfrac{2 \sin^2 \dfrac{x}{2}}{x} = 0
	\]
	所以
	\[
	1 - \cos x = o(x) \quad (x \to 0)
	\]
	
	(2)
	\[
	\lim_{x \to 0} \frac{\tan x - \sin x}{x^2} = \lim_{x \to 0} \left( \frac{\sin x}{x \cos x} \cdot \frac{1 - \cos x}{x} \right) = 0
	\]
	所以
	\[
	\tan x - \sin x = o(x^2) \quad (x \to 0)
	\]
	
	(3)$x \to 0 $时，$x \sin \dfrac{1}{x}$ 是无穷小量，且$\left| \dfrac{x \sin \dfrac{1}{x}}{x} \right| \leqslant 1$，因此 
	\[
	x \sin \frac{1}{x} = O(x) \quad (x \to 0)
	\]
	
	(4)
	\[
	\lim_{x \to 0} \frac{\sin x}{x} = 1
	\]
	所以
	\[
	\sin x \sim x \quad (x \to 0) \quad \text{ 或者 } \quad \sin x = x + o(x) \quad (x \to 0)
	\]
	
	(5)
	\[
	\lim_{x \to 0} \dfrac{1 - \cos x}{\dfrac{1}{2} x^2} = \lim_{x \to 0} \dfrac{2 \sin^2 \dfrac{x}{2}}{\dfrac{1}{2} x^2} = 1
	\]
	所以
	\[
	1 - \cos x \sim \frac{1}{2} x^2 \quad (x \to 0) \quad \text{ 或者 } \quad 1 - \cos x = \frac{1}{2} x^2 + o(x^2) \quad (x \to 0)
	\]
	
	(6)
	\[
	\lim_{x \to 0} \frac{\tan x - \sin x}{\dfrac{1}{2} x^3} = \lim_{x \to 0} \left( \frac{\sin x}{x \cos x} \cdot \dfrac{1 - \cos x}{\dfrac{1}{2} x^2} \right) = 1
	\]
	所以
	\[
	\tan x - \sin x \sim \dfrac{1}{2} x^3 \quad (x \to 0)\text{或者}\tan x - \sin x=\frac{1}{2}x^3+o(x^3)\quad(x\to0)\]
\end{example}
我们往往取 $v(x) = (x - x_0)^k$ 与 $u(x)$ 进行比较 ($x \to x_0$)，若 $x \to \infty$，
则取 $v(x) = \dfrac{1}{x^k}$.我们称 $k$ 为 $u(x)$ 作为无穷小量的阶数.例如
\[ 1 - \cos x \sim \frac{1}{2}x^2 \quad (x \to 0), \]
得 $x \to 0$时，$1 - \cos x$ 是二阶无穷小量.

类似的，对无穷大量我们有类似的定义
\begin{definition}
	对于 $u(x), v(x)$，在 $x \to x_0$ 时，是无穷大量，那么：
	
	(1) 若 $\lim\limits_{x \to x_0} \dfrac{u(x)}{v(x)} = \infty$，则称 $x \to x_0$ 时，$u(x)$ 关于 $v(x)$ 是高阶无穷大量，或 $v(x)$ 关于 $u(x)$ 是低阶无穷大量.
	记作\[v(x)=o(u(x))\]
	(2) 若 $\exists A > 0$，当 $x$ 在 $x_0$ 的某去心邻域中，成立
	\[ \left| \frac{u(x)}{v(x)} \right| \leqslant A \]
	则称 $x \to x_0$ 时，$\dfrac{u(x)}{v(x)}$ 是有界量，记作
	\[ u(x) = O(v(x)) \quad (x \to x_0). \]
	
	(3) 若 $\exists A, a > 0$，当 $x$ 在 $x_0$ 的某去心邻域中，成立
	\[ a \leqslant \left| \frac{u(x)}{v(x)} \right| \leqslant A \]
	则称 $x \to x_0$ 时，$u(x)$ 与 $v(x)$ 是同阶无穷大量.
	
	若$\lim\limits_{x \to x_0} \dfrac{u(x)}{v(x)} = C \neq 0$，则 $u(x)$ 与 $v(x)$ 是同阶无穷大量.
	
	(4) 若 $\lim\limits_{x \to x_0} \dfrac{u(x)}{v(x)} = 1$，则称 $x \to x_0$ 时，$u(x)$ 与 $v(x)$ 是等价无穷大量，记作
	\[ u(x) \sim v(x) \quad (x \to x_0). \]
	
\end{definition}
\begin{example}
	
	(1) $\lim\limits_{x \to +\infty} \dfrac{a^k}{x^k} = \infty \quad (a > 1) \Rightarrow a^k \text{ 关于 } x^k\text{ 是高阶无穷大量.}$
	
	(2) $\lim\limits_{x \to +\infty} \dfrac{\ln^k x}{x} = 0 \Rightarrow \ln^k x \text{ 关于 } x \text{ 是低阶无穷大量.}$
	
	(3) $x \to +\infty$ 时， $x(\arctan x + \sin x)$ 与 $x$ 是无穷大量，且
	\[ \left| \frac{x(\arctan x + \sin x)}{x} \right| \leqslant 3 \]
	从而
	\[ x(\arctan x + \sin x) = O(x) \quad (x \to +\infty) \]
	
	(4)$$\lim_{x \to +\infty} \dfrac{x^3 \sin \dfrac{1}{x}}{x^2} = \lim_{x \to +\infty} \dfrac{\sin \dfrac{1}{x}}{\dfrac{1}{x}} = 1 $$
	所以$$x^3 \sin \frac{1}{x} \sim x^2 \quad (x \to +\infty)$$
	
\end{example}
下面是一些常用的等价关系，他们常常在极限的运算中被使用
\begin{proposition}
	\( x \to 0 \) 时
	
	(1) \(\ln(1+x) \sim x\)
	
	(2) \(e^x - 1 \sim x\)
	
	(3) \((1+x)^\alpha - 1 \sim \alpha x\)
	
	(4) \(\sin x \sim x\), \(1 - \cos x \sim \dfrac{1}{2}x^2\), \(\tan x \sim x\), \(\tan x - \sin x \sim \dfrac{1}{2}x^3\)
\end{proposition}
\begin{proof}
	(1) 由 \(\lim\limits_{x \to \infty} (1+\dfrac{1}{x})^{x} = e\)，得\(\lim\limits_{x \to 0}(1+x)^{\frac{1}{x}}  = e\)
	
	取对数
	\[
	\lim_{x \to 0} \frac{\ln(1+x)}{x} = \lim_{x \to 0} \ln(1+x)^{\frac{1}{x}} = 1
	\]
	
	(2) 设 \( y = e^x - 1 \)，当 \( x \to 0 \), \( y \to 0 \), 且 \( x = \ln(1+y) \)，则
	\[
	\lim_{x \to 0} \frac{e^x - 1}{x} = \lim_{y \to 0} \frac{y}{\ln(1+y)} = 1
	\]
	
	(3) 令 \( y = (1+x)^\alpha - 1 \)，当 \( x \to 0 \) 时，\( y \to 0 \)，
	\[
	\lim_{x \to 0} \frac{(1+x)^\alpha - 1}{x} = \lim_{x \to 0} \frac{(1+x)^\alpha - 1}{\ln(1+x)^\alpha} \cdot \frac{\alpha \ln(1+x)}{x} = \lim_{y \to 0} \frac{y}{\ln(1+y)} \cdot \lim_{x \to 0} \frac{\alpha \ln(1+x)}{x} = \alpha
	\]
\end{proof}
\begin{proposition}
	设 \( u(x), v(x) \) 和 \( w(x) \) 在 \( x_0 \) 的某个去心邻域内有定义，
	且 \( v(x) \sim w(x) \quad (x \to x_0) \)，那么：
	
	(1) 若 \(\lim\limits_{x \to x_0} u(x) w(x) = A\)，则 \(\lim\limits_{x \to x_0} u(x) v(x) = A\).
	
	(2) 若 \(\lim\limits_{x \to x_0} \dfrac{u(x)}{w(x)} = A\)，则 \(\lim\limits_{x \to x_0} \dfrac{u(x)}{v(x)} = A\).
	
	
	
\end{proposition}
\begin{proof}
	由极限的运算法则得证.
\end{proof}
下面是一些例子

\begin{example}
	(1)\[	\lim_{x \to 0} (\cos x)^{1/x^2} = \lim_{x \to 0} \left[1 - (1 - \cos x)\right]^{1/x^2} = \lim_{x \to 0} \left(1 - \frac{x^2}{2}-o(x^2)\right)^{1/x^2} = \dfrac{1}{\sqrt{e}}\]
	(2)\[	\lim_{x \to \infty} \frac{a_n x^n + a_{n+1} x^{n+1} + \cdots + a_m x^m}{b_n x^n + b_{n+1} x^{n+1} + \cdots + b_m x^m} = \lim_{x \to \infty} \frac{a_m x^m}{b_m x^m} = \frac{a_m}{b_m} \quad (a_m, b_m \neq 0)\]
	\[	\lim_{x \to 0} \frac{A_n x^n + A_{n+1} x^{n+1} + \cdots + A_m x^m}{B_n x^n + B_{n+1} x^{n+1} + \cdots + B_m x^m} = \lim_{x \to 0} \frac{A_n x^n}{B_n x^n} = \frac{A_n}{B_n} \quad (A_n, B_n \neq 0)	\]
\end{example}
\begin{example}\label{阶的比较}
	当 \( x \to +\infty \) 时，下列无穷大量从低阶到高阶的排列次序：
	\[ \ln^k x \quad (k > 0), \quad x^\alpha \quad (\alpha > 0), \quad a^x \quad (a > 1), \quad [x]!, \quad x^x \]
	
	
\end{example}
\begin{proof}
	设 \( n \leqslant x < n+1 \)，则
	\[ 0 < \frac{x^\alpha}{a^x} < \frac{(n+1)^\alpha}{a^n}, \quad 0 < \frac{a^x}{[x]!} < \frac{a^{n+1}}{n!}, \quad 0 < \frac{[x]!}{x^x} < \frac{(n+1)!}{n^n} \]
	
	由(\ref{几个特殊的序列})得\(\lim\limits_{n \to \infty} \dfrac{(n+1)^\alpha}{a^n} = 0\)
	
	对 \(\lim\limits_{n \to \infty} \dfrac{a^{n+1}}{n!} = 0\)，考虑 \(\lim\limits_{n \to \infty} \dfrac{a^n}{n!}\).
	
	对 \(n \ge 2a - 1\)，
	\[ 0 < \frac{a^n}{n!} < \frac{a^{2a-1}}{(2a-1)!} \left(\frac{1}{2}\right)^{n-2a+1} \]
	右端趋于 0.
	
	对 \(\lim\limits_{n \to \infty} \dfrac{(n+1)!}{n^n} = 0\)，设 \(a_n = \dfrac{(n+1)!}{n^n}\)，则
	\[
	\dfrac{a_{n+1}}{a_n} = \dfrac{\dfrac{(n+2)!}{(n+1)^{n+1}}}{\dfrac{(n+1)!}{n^n}} = \left(1 - \frac{1}{n+1}\right)^{n-1} \cdot \left(1 - \frac{1}{(n+1)^2}\right)<1
	\]
	
	
	\(\{a_n\}\) 递减有下界，则$\{a_n\}$的极限存在，且
	\[a_{n+1}=(\frac{n+2}{n+1})(1-\frac{1}{n+1})^na_n\]
	两侧取$n\to\infty$
	\[
	\lim_{n \to \infty} a_n = 0
	\]
	
	得
	\[
	\lim_{x \to +\infty} \frac{x^\alpha}{a^x} = 0, \quad \lim_{x \to +\infty} \frac{a^x}{x!} = 0, \quad \lim_{x \to +\infty} \frac{[x]!}{x^x} = 0
	\]
	
	同时
	\[
	\lim_{x \to +\infty} \frac{\ln^k x}{x^\alpha} = \lim_{y \to +\infty} \frac{y^k}{(e^y)^\alpha} = 0 \quad (y = \ln x)\]
\end{proof}
